\section{Stratified Optical Model}
\label{sec:model}

\subsection{Geometry and assumptions}

We consider a two-dimensional vertical section through a horizontally
homogeneous medium. Let $x$ denote horizontal distance and let $h\geq 0$
denote height above an opaque planar boundary located at $h=0$. The
refractive index is assumed to be isotropic and vertically stratified,
\begin{equation}
    n=n(h),
    \qquad n(h)>0,
\end{equation}
with
\begin{equation}
    n_0\equiv n(0).
\end{equation}
Throughout the paper, $h$ is measured upward from the opaque boundary.
Horizontal translational symmetry allows each ray to be analysed in a single
vertical plane. A ray is represented by a curve
$h=h(x)$ wherever it is single-valued in $x$.

The boundary at $h=0$ is assumed to be perfectly opaque. A source and an
observer are mutually visible only if at least one geometrical-optics ray
connecting them remains in the open half-space $h>0$, except possibly at a
single tangency point in the limiting case.

\subsection{Fermat's principle and the first integral}

The optical path length of a ray between two fixed endpoints is, in the
geometrical-optics approximation \cite{born1999principles},
\begin{equation}
    \mathcal{S}[h]
    =
    \int n(h)\,ds
    =
    \int n(h)\sqrt{1+h'^2}\,dx,
    \label{eq:fermat-functional}
\end{equation}
where $h'=dh/dx$. Geometrical-optics rays are stationary paths of
the functional in Eq.~\eqref{eq:fermat-functional}. The corresponding
Lagrangian is
\begin{equation}
    \mathcal{L}(h,h')=n(h)\sqrt{1+h'^2}.
\end{equation}
It contains no explicit dependence on $x$. The associated horizontal
translational symmetry therefore supplies a first integral: the Beltrami
identity reduces the Euler--Lagrange equation to
\begin{equation}
    \mathcal{L}
    -h'\frac{\partial\mathcal{L}}{\partial h'}
    =
    \frac{n(h)}{\sqrt{1+h'^2}}
    =a,
    \label{eq:first-integral}
\end{equation}
where $a$ is constant along a ray. Introducing the angle $\theta$ between
ray and horizontal, so that $\cos\theta=1/\sqrt{1+h'^2}$,
Eq.~\eqref{eq:first-integral} becomes
\begin{equation}
    n(h)\cos\theta=a.
\end{equation}
This is the continuous form of Snell's law appropriate to a horizontally
stratified medium \cite{schroeder2000astronomical,nener2003analytical}.

Rearranging Eq.~\eqref{eq:first-integral} gives the ray slope in the form
\begin{equation}
    h'^2
    =
    \frac{n^2(h)}{a^2}-1.
    \label{eq:ray-slope}
\end{equation}
On any monotonic branch, inverting this relation gives
\begin{equation}
    \left|\frac{dx}{dh}\right|
    =
    \frac{a}{\sqrt{n^2(h)-a^2}}.
    \label{eq:horizontal-increment-general}
\end{equation}
These relations support both the direct construction of ray trajectories and
the inverse reconstruction of a refractive-index profile from a prescribed
visibility law.

\subsection{Grazing rays and horizon distance}

For fixed endpoint heights, the limiting visible ray just grazes the opaque
boundary. At its point of contact with $h=0$ the ray is
horizontal, so that $h'=0$. Equation~\eqref{eq:first-integral} then fixes its
constant to
\begin{equation}
    a=n_0.
    \label{eq:grazing-constant}
\end{equation}
Such a grazing ray requires $n(h)\geq n_0$ throughout the heights it
traverses. Hence the upward-bending profiles considered here satisfy
$n(h)\geq n_0$ near the boundary.

Define $d(h)$ as the maximum horizontal distance from a boundary tangency
point to an endpoint at height $h$. Using Eq.~\eqref{eq:grazing-constant} in
Eq.~\eqref{eq:horizontal-increment-general} gives
\begin{equation}
    d(h)
    =
    \int_0^h
    \frac{n_0\,d\eta}
         {\sqrt{n^2(\eta)-n_0^2}}.
    \label{eq:horizon-distance-general}
\end{equation}
The square-root singularity at the lower limit is integrable whenever the
profile rises regularly from $n_0$.

Let an observer and a target have heights $h_{\mathrm{o}}$ and
$h_{\mathrm{t}}$, respectively. By horizontal translational symmetry, the
grazing point may be chosen as the origin. The limiting ray then consists
of two reciprocal branches meeting tangentially at the boundary, and the
maximum horizontal separation is
\begin{equation}
    L_{\max}(h_{\mathrm{o}},h_{\mathrm{t}})
    =
    d(h_{\mathrm{o}})+d(h_{\mathrm{t}}).
    \label{eq:maximum-separation-general}
\end{equation}
Figure~\ref{fig:grazing-geometry} summarizes the corresponding one-sided
horizon distance and the two-sided limiting geometry.
\begin{figure}[tb]
    \centering
    \resizebox{\textwidth}{!}{% Figure 1: grazing-ray geometry.
% This file is included from sections/mathematical_model.tex.
\begin{tikzpicture}[
    font=\small,
    ray/.style={draw=black, line width=0.95pt, -{Stealth[length=2.0mm,width=1.3mm]}},
    boundary/.style={draw=black, line width=1.25pt},
    construction/.style={draw=black!55, densely dashed, line width=0.45pt},
    dimension/.style={{Stealth[length=1.8mm,width=1.2mm]}-{Stealth[length=1.8mm,width=1.2mm]}, draw=black!75, line width=0.5pt},
    tangent/.style={draw=black!65, line width=0.55pt},
    point/.style={circle, fill=black, inner sep=1.45pt},
    paneltitle/.style={font=\small},
    label/.style={font=\small}
]

% Common dimensions in drawing units.
\def\panelW{6.0}
\def\baseY{0}
\def\obsX{0.75}
\def\grazeX{5.05}
\def\obsH{2.15}

% --------------------------------------------------------------------------
% (a) Single-point horizon distance
% --------------------------------------------------------------------------
\begin{scope}[xshift=0cm]
  \node[paneltitle, anchor=west] at (0.0,3.05) {(a) Single-point horizon distance};

  % Opaque boundary.
  \draw[boundary] (0,\baseY) -- (\panelW,\baseY);
  \node[anchor=west] at (0.02,-0.36) {$h=0$};

  % Exact exponential-model grazing branch, normalized and rescaled.
  % Parameter u runs from 0 at G to a positive value at O, with
  % h/R = -ln(cos u).  The affine scaling only sets the panel proportions.
  \draw[ray, samples=90, smooth, variable=\u, domain=0:1.08]
    plot ({\grazeX-(\grazeX-\obsX)*(\u/1.08)},
          {2.15*(-ln(cos(\u r)))/(-ln(cos(1.08 r))) });

  \coordinate (Oa) at (\obsX,\obsH);
  \coordinate (Ga) at (\grazeX,0);
  \node[point] at (Oa) {};
  \node[point] at (Ga) {};
  \node[anchor=south east] at ($(Oa)+(-0.05,0.08)$) {$O$};
  \node[anchor=north] at ($(Ga)+(0,-0.10)$) {$G$};

  % Height and horizontal-distance constructions.
  \draw[construction] (Oa) -- (\obsX,0);
  \draw[dimension] (\obsX-0.22,0.02) -- node[left] {$h$} (\obsX-0.22,\obsH-0.02);
  \draw[dimension] (\obsX,-0.63) -- node[fill=white,inner sep=1.2pt] {$d(h)$} (\grazeX,-0.63);

  % Tangent at grazing point.
  \draw[tangent] ($(Ga)+(-0.38,0.10)$) -- ($(Ga)+(0.38,0.10)$);
  \node[anchor=south] at (3.20,1.10) {grazing ray};
\end{scope}

% --------------------------------------------------------------------------
% (b) Maximum mutual-visibility distance
% --------------------------------------------------------------------------
\begin{scope}[xshift=7.15cm]
  \node[paneltitle, anchor=west] at (0.0,3.05) {(b) Maximum mutual-visibility distance};

  \def\leftX{0.75}
  \def\gX{3.20}
  \def\rightX{5.45}
  \def\leftH{2.12}
  \def\rightH{1.58}

  \draw[boundary] (0,0) -- (\panelW,0);
  \node[anchor=west] at (0.02,-0.36) {$h=0$};

  % Left exact grazing branch.
  \draw[ray, samples=75, smooth, variable=\u, domain=0:1.08]
    plot ({\gX-(\gX-\leftX)*(\u/1.08)},
          {\leftH*(-ln(cos(\u r)))/(-ln(cos(1.08 r))) });
  % Right reciprocal branch; arrow points toward the target.
  \draw[ray, samples=75, smooth, variable=\u, domain=0:0.96]
    plot ({\gX+(\rightX-\gX)*(\u/0.96)},
          {\rightH*(-ln(cos(\u r)))/(-ln(cos(0.96 r))) });

  \coordinate (Ob) at (\leftX,\leftH);
  \coordinate (Gb) at (\gX,0);
  \coordinate (Tb) at (\rightX,\rightH);
  \node[point] at (Ob) {};
  \node[point] at (Gb) {};
  \node[point] at (Tb) {};
  \node[anchor=south east] at ($(Ob)+(-0.04,0.08)$) {$O$};
  \node[anchor=north] at ($(Gb)+(0,-0.10)$) {$G$};
  \node[anchor=south west] at ($(Tb)+(0.06,0.06)$) {$T$};

  \draw[construction] (Ob) -- (\leftX,0);
  \draw[construction] (Tb) -- (\rightX,0);
  \draw[dimension] (\leftX-0.22,0.02) -- node[left] {$h_{\mathrm{o}}$} (\leftX-0.22,\leftH-0.02);
  \draw[dimension] (\rightX+0.22,0.02) -- node[right] {$h_{\mathrm{t}}$} (\rightX+0.22,\rightH-0.02);

  \draw[tangent] ($(Gb)+(-0.38,0.10)$) -- ($(Gb)+(0.38,0.10)$);
  \node[anchor=south] at (3.18,1.10) {grazing ray};

  \draw[dimension] (\leftX,-0.60) -- node[fill=white,inner sep=1.0pt] {$d(h_{\mathrm{o}})$} (\gX,-0.60);
  \draw[dimension] (\gX,-0.60) -- node[fill=white,inner sep=1.0pt] {$d(h_{\mathrm{t}})$} (\rightX,-0.60);
  \draw[dimension] (\leftX,-1.02) -- node[fill=white,inner sep=1.0pt]
      {$L_{\max}=d(h_{\mathrm{o}})+d(h_{\mathrm{t}})$} (\rightX,-1.02);
\end{scope}

\end{tikzpicture}
}
    \caption{Grazing-ray geometry above an opaque planar boundary.
    (a) The horizon distance $d(h)$ from a boundary tangency point $G$ to an
    endpoint $O$ at height $h$. (b) The limiting mutual-visibility ray between
    an observer at height $h_{\mathrm{o}}$ and a target at height
    $h_{\mathrm{t}}$; its two reciprocal branches meet with a horizontal
    tangent at $G$, giving
    $L_{\max}=d(h_{\mathrm{o}})+d(h_{\mathrm{t}})$.}
    \label{fig:grazing-geometry}
\end{figure}
For separations below this limit, an unobstructed ray exists. At equality,
the limiting ray grazes the boundary. Beyond it, every connecting ray
intersects the opaque surface.
Equation~\eqref{eq:horizon-distance-general} therefore provides the basic
link between a stratified refractive-index profile and the resulting
horizon-occlusion law.
